%! Lengths and Angles from Dot Products — Unit 1, Lesson 1.2 — Exit Ticket (Answer Key)
\documentclass[10pt]{article}
\usepackage{linalg-article}
\usepackage{linalg-key}   % pulls in -boxes; do NOT also load -boxes
\newcommand{\TallMath}[1]{$\displaystyle #1\rule[-1.4em]{0pt}{3.2em}$}
\newcommand{\vv}{\mathbf{v}}
\newcommand{\ww}{\mathbf{w}}

\begin{document}

\pageheader{Unit 1, Lesson 1.2}{Exit Ticket --- Answer Key}
\namedateperiod

\vspace{0.15in}

Show your work.

\begin{enumerate}[leftmargin=1.8em, itemsep=14pt]
  \item Compute the dot product $\begin{bmatrix} 3 \\ 4 \end{bmatrix}\cdot\begin{bmatrix} 1 \\ 2 \end{bmatrix}$.
        \ansline{$(3)(1)+(4)(2)=3+8=11$}
  \item Find the length of $\begin{bmatrix} 6 \\ 8 \end{bmatrix}$ using $\|\vv\|=\sqrt{\vv\cdot\vv}$.
        \ansline{$\sqrt{6^2+8^2}=\sqrt{36+64}=\sqrt{100}=10$}
  \item Are $\begin{bmatrix} 4 \\ 1 \end{bmatrix}$ and $\begin{bmatrix} -1 \\ 4 \end{bmatrix}$
        perpendicular? Use the dot product to decide, and say what your answer means about the
        angle between them.
        \ansline{$(4)(-1)+(1)(4)=-4+4=0$. The dot product is $0$, so yes --- they are
        perpendicular: the angle between them is $90^\circ$.}
\end{enumerate}

\begin{teachernote}
Item 3 checks the central idea --- dot product $=0$ means a right angle. Full credit needs both
the computation ($=0$) \emph{and} the meaning (perpendicular / $90^\circ$). A common slip in
item~1 is crossing components ($3\cdot2+4\cdot1$); look for matching products $3\cdot1+4\cdot2$.
\end{teachernote}

\end{document}
